\renewcommand{\thesection}{\arabic{section}}\renewcommand{\theHsection}{chapter.\arabic{section}}\setcounter{section}{2}\section{单变量函数的微分学}
\noindent\begin{minipage}{\textwidth}\ProblemHeading{例题}{3.1}
设 $y=c$，其中 $c$ 为常数。求 $y'$。
\end{minipage}\par\vfill
\noindent\begin{minipage}{\textwidth}\ProblemHeading{例题}{3.2}
设 $y=x^n$，$x\in\mathbb R$，其中 $n$ 为正整数。求 $y'$。
\end{minipage}\par\vfill
\noindent\begin{minipage}{\textwidth}\ProblemHeading{例题}{3.3}
求正弦函数和余弦函数的导函数。
\end{minipage}\par\vfill
\noindent\begin{minipage}{\textwidth}\ProblemHeading{例题}{3.4}
设 $a>0$ 且 $a\ne1$。求对数函数 $y=\log_a x\ (x>0)$ 的导函数。
\end{minipage}\par\vfill
\newpage
\noindent\begin{minipage}{\textwidth}\ProblemHeading{例题}{3.5}
求分段函数\[f(x)=\begin{cases}x^3,&x\ge0,\\x^2,&x<0\end{cases}\]的导函数。
\end{minipage}\par\vfill
\noindent\begin{minipage}{\textwidth}\ProblemHeading{例题}{3.6（改写）}
证明函数 $f(x)=|x|$ 在 $x=0$ 处连续，但不可导。
\end{minipage}\par\vfill
\noindent\begin{minipage}{\textwidth}\ProblemHeading{例题}{3.7（改写）}
证明函数 $f(x)=\sqrt[3]x$ 在 $x=0$ 处连续，但不可导。
\end{minipage}\par\vfill
\noindent\begin{minipage}{\textwidth}\ProblemHeading{例题}{3.8（改写）}
设 $f(x)$ 可导，$c$ 为常数。证明 $(cf(x))'=cf'(x)$。
\end{minipage}\par\vfill
\newpage
\noindent\begin{minipage}{\textwidth}\ProblemHeading{例题}{3.9}
求 $f(x)=x^2(\sin x+\cos x)$ 的导数。
\end{minipage}\par\vfill
\noindent\begin{minipage}{\textwidth}\ProblemHeading{例题}{3.10}
求函数 $\tan x$ 与 $\cot x$ 的导数。
\end{minipage}\par\vfill
\noindent\begin{minipage}{\textwidth}\ProblemHeading{例题}{3.11}
求函数 $\sec x$ 与 $\csc x$ 的导数。
\end{minipage}\par\vfill
\noindent\begin{minipage}{\textwidth}\ProblemHeading{例题}{3.12}
求 $\sin(x^3)$ 的导数。
\end{minipage}\par\vfill
\newpage
\noindent\begin{minipage}{\textwidth}\ProblemHeading{例题}{3.13}
求 $z=\sin(\cos(x^2))$ 的导数。
\end{minipage}\par\vfill
\noindent\begin{minipage}{\textwidth}\ProblemHeading{例题}{3.14}
求 $z=(1-x)^9$ 的导数。
\end{minipage}\par\vfill
\noindent\begin{minipage}{\textwidth}\ProblemHeading{例题}{3.15}
求 $\displaystyle z=\left(\frac{1+x}{1-x}\right)^3$ 的导数。
\end{minipage}\par\vfill
\noindent\begin{minipage}{\textwidth}\ProblemHeading{例题}{3.16}
设 $f$ 在点 $x$ 处可导且 $f(x)\ne0$。证明 $\ln|f|$ 在该点可导，且\[(\ln|f|)'(x)=\frac{f'(x)}{f(x)}.\]特别地，证明 $(\ln|x|)'=1/x\ (x\ne0)$。
\end{minipage}\par\vfill
\newpage
\noindent\begin{minipage}{\textwidth}\ProblemHeading{例题}{3.17}
求下列反三角函数在其可导区间内的导函数。
\SubQ{1}{$y=\arcsin x$。}
\SubQ{2}{$y=\arccos x$。}
\SubQ{3}{$y=\arctan x$。}
\end{minipage}\par\vfill
\noindent\begin{minipage}{\textwidth}\ProblemHeading{例题}{3.18}
求下列反三角函数在其可导区间内的导函数。
\SubQ{1}{$y=\operatorname{arccot}x$，其值域取 $(0,\pi)$。}
\end{minipage}\par\vfill
\noindent\begin{minipage}{\textwidth}\ProblemHeading{例题}{3.19}
求指数函数 $y=a^x$ 的导函数，其中 $a>0$ 且 $a\ne1$。
\end{minipage}\par\vfill
\noindent\begin{minipage}{\textwidth}\ProblemHeading{例题}{3.20}
设 $y=x^\alpha\ (x>0)$，其中 $\alpha$ 为任意实数。证明 $(x^\alpha)'=\alpha x^{\alpha-1}$。
\end{minipage}\par\vfill
\newpage
\noindent\begin{minipage}{\textwidth}\ProblemHeading{例题}{3.21（改写）}
设 $u(x),v(x)$ 可导且 $v(x)>0$。证明 $y=v(x)^{u(x)}$ 可导，并求其导数。
\end{minipage}\par\vfill
\noindent\begin{minipage}{\textwidth}\ProblemHeading{例题}{3.22（校订）}
分别在单调区间 $x<-1$ 与 $x>-1$ 上，求 $y=xe^x$ 的反函数 $x=x(y)$ 的导数，结果可用 $x$ 表示。
\end{minipage}\par\vfill
\noindent\begin{minipage}{\textwidth}\ProblemHeading{例题}{3.23}
设 $P_n(x)=a_nx^n+a_{n-1}x^{n-1}+\cdots+a_1x+a_0$ 是 $n$ 次多项式。
\SubQ{1}{求 $P_n$ 的各阶导数。}
\SubQ{2}{若 $x_0$ 为 $P_n$ 的 $r$ 重根，即 $P_n(x)=(x-x_0)^rQ_{n-r}(x)$ 且 $Q_{n-r}(x_0)\ne0$，证明 $P_n(x_0)=P_n^{\prime}(x_0)=\cdots=P_n^{(r-1)}(x_0)=0$，而 $P_n^{(r)}(x_0)\ne0$。}
\end{minipage}\par\vfill
\noindent\begin{minipage}{\textwidth}\ProblemHeading{例题}{3.24}
求 $f(x)=e^{ax}\ (x\in\mathbb R)$ 的各阶导数，其中 $a$ 为常数。
\end{minipage}\par\vfill
\newpage
\noindent\begin{minipage}{\textwidth}\ProblemHeading{例题}{3.25}
求 $f(x)=\ln(1+x)\ (x>-1)$ 的 $n$ 阶导数，$n\ge1$。
\end{minipage}\par\vfill
\noindent\begin{minipage}{\textwidth}\ProblemHeading{例题}{3.26}
求 $\sin x$ 与 $\cos x$ 的 $n$ 阶导数，$x\in\mathbb R$，$n\ge1$。
\end{minipage}\par\vfill
\noindent\begin{minipage}{\textwidth}\ProblemHeading{例题}{3.27}
求 $(1+x)^\alpha\ (x>-1)$ 的 $n$ 阶导数，其中 $\alpha\in\mathbb R$，$n\ge1$。
\end{minipage}\par\vfill
\noindent\begin{minipage}{\textwidth}\ProblemHeading{例题}{3.28}
设 $y=\arctan x$，求 $y^{(n)}(0)\ (n\ge1)$。
\end{minipage}\par\vfill
\newpage
\noindent\begin{minipage}{\textwidth}\ProblemHeading{习题}{3.29}
讨论下列函数在 $x=0$ 处是否可导。
\SubQ{1}{$f(x)=|\sin x|$。}
\SubQ{2}{$f(x)=\begin{cases}x+1,&x\ge0,\\1,&x<0.\end{cases}$}
\SubQ{3}{$f(x)=\begin{cases}x^2\sin(1/x),&x\ne0,\\1,&x=0.\end{cases}$}
\end{minipage}\par\vfill
\noindent\begin{minipage}{\textwidth}\ProblemHeading{习题}{3.30}
讨论下列函数在 $x=0$ 处是否可导。
\SubQ{1}{$f(x)=\begin{cases}\ln(1+x),&x\ge0,\\x+1,&x<0.\end{cases}$}
\SubQ{2}{$f(x)=|x|e^x$。}
\SubQ{3}{$f(x)=|x^3|$。}
\end{minipage}\par\vfill
\noindent\begin{minipage}{\textwidth}\ProblemHeading{习题}{3.31}
求 $a,b$ 的值，使下列函数在其定义域内处处可导。
\SubQ{1}{$f(x)=\begin{cases}x^2,&x\le1,\\ax+b,&x>1.\end{cases}$}
\SubQ{2}{$f(x)=\begin{cases}\ln(1+x),&-1<x<0,\\ax+b,&x\ge0.\end{cases}$}
\end{minipage}\par\vfill
\noindent\begin{minipage}{\textwidth}\ProblemHeading{习题}{3.32}
设 $g(x)$ 在 $x=a$ 处连续，记 $f(x)=(x-a)g(x)$。证明 $f'(a)=g(a)$。
\end{minipage}\par\vfill
\newpage
\noindent\begin{minipage}{\textwidth}\ProblemHeading{习题}{3.33}
设 $f(x)$ 在 $x_0$ 处可导，$\alpha,\beta$ 为常数。证明\[\lim_{h\to0}\frac{f(x_0+\alpha h)-f(x_0-\beta h)}h=(\alpha+\beta)f'(x_0).\]
\end{minipage}\par\vfill
\noindent\begin{minipage}{\textwidth}\ProblemHeading{习题}{3.34}
设 $f(x)$ 在 $x=a$ 处可导且 $f(a)\ne0$。证明 $|f(x)|$ 在 $x=a$ 处也可导。若 $f(a)=0$，结论是否仍成立？
\end{minipage}\par\vfill
\noindent\begin{minipage}{\textwidth}\ProblemHeading{习题}{3.35}
求下列函数的导数。
\SubQ{1}{$\displaystyle y=\frac{3x^2+9x-2}{5x+8}$。}
\SubQ{2}{$y=\sin x\tan x+\cot x$。}
\SubQ{3}{$y=x^2\log_3x$。}
\end{minipage}\par\vfill
\noindent\begin{minipage}{\textwidth}\ProblemHeading{习题}{3.36}
求下列函数的导数。
\SubQ{1}{$\displaystyle y=\frac{x}{1-\cos x}$。}
\SubQ{2}{$\displaystyle y=\frac{1+\ln x}{1-\ln x}$。}
\SubQ{3}{$\displaystyle y=\frac{(1+x^2)\ln x}{\sin x+\cos x}$。}
\end{minipage}\par\vfill
\newpage
\noindent\begin{minipage}{\textwidth}\ProblemHeading{习题}{3.37}
求下列函数的导数。
\SubQ{1}{$y=(x^2+1)(3x-1)(1-x^3)$。}
\SubQ{2}{$y=x^3\tan x\ln x$。}
\end{minipage}\par\vfill
\noindent\begin{minipage}{\textwidth}\ProblemHeading{习题}{3.38}
求下列函数的导数。
\SubQ{1}{$y=x\sqrt{1-x^2}$。}
\SubQ{2}{$y=\sqrt[3]{1+\ln^2x}$。}
\SubQ{3}{$\displaystyle y=\arccos\frac{2x-1}{\sqrt3}$。}
\end{minipage}\par\vfill
\noindent\begin{minipage}{\textwidth}\ProblemHeading{习题}{3.39}
求下列函数的导数。
\SubQ{1}{$y=(\sin x+\cos x)^3$。}
\SubQ{2}{$y=(\sin(x^3))^3$。}
\SubQ{3}{$y=\sqrt{x+\sqrt{x+\sqrt x}}$。}
\end{minipage}\par\vfill
\noindent\begin{minipage}{\textwidth}\ProblemHeading{习题}{3.40}
求下列函数的导数。
\SubQ{1}{$y=\sin(\sin(\sin x))$。}
\SubQ{2}{$y=\sin(\cos^5(\arctan(x^3)))$。}
\SubQ{3}{$\displaystyle y=\left(\frac{x^3-1}{x^4+1}\right)^3$。}
\end{minipage}\par\vfill
\newpage
\noindent\begin{minipage}{\textwidth}\ProblemHeading{习题}{3.41}
求下列函数的导数。
\SubQ{1}{$y=x\sqrt{1+x^2}\sin x$。}
\SubQ{2}{$y=e^{\sqrt{x^2+1}}$。}
\SubQ{3}{$y=\ln\bigl(\ln^2(\ln^3x)\bigr)$。}
\end{minipage}\par\vfill
\noindent\begin{minipage}{\textwidth}\ProblemHeading{习题}{3.42}
求下列函数的导数。
\SubQ{1}{$y=x^{x^x}+x^x+x^{2^x}$。}
\SubQ{2}{$y=(\ln x)^{e^x}$。}
\SubQ{3}{$y=(\tan x)^{\cot x}$。}
\end{minipage}\par\vfill
\noindent\begin{minipage}{\textwidth}\ProblemHeading{习题}{3.43}
求下列函数的导数。
\SubQ{1}{$y=10^x(\sin x)^{\cos x}$。}
\SubQ{2}{$\displaystyle y=\frac{(x+5)^2(x-4)^{1/3}}{(x+2)^5(x+4)^{1/2}}$。}
\SubQ{3}{$\displaystyle y=\frac{x^2}{1-x}\sqrt{\frac{x+1}{1+x+x^2}}$。}
\end{minipage}\par\vfill
\noindent\begin{minipage}{\textwidth}\ProblemHeading{习题}{3.44}
设 $f(x)=x^3$，求 $f'(x^2)$ 与 $(f(x^2))'$。
\end{minipage}\par\vfill
\newpage
\noindent\begin{minipage}{\textwidth}\ProblemHeading{习题}{3.45}
设 $f(x)=\ln(x+\sqrt{1+x^2})$，$g(x)=e^{\sqrt{x^2+1}}$。求 $f'(g(x))$ 与 $(f(g(x)))'$。
\end{minipage}\par\vfill
\noindent\begin{minipage}{\textwidth}\ProblemHeading{习题}{3.46}
设 $f(x)$ 处处可导，求下列各式的 $\dfrac{dy}{dx}$。
\SubQ{1}{$y=f(x^3)$。}
\SubQ{2}{$y=f(\sin^2x)+f(\cos^2x)$。}
\SubQ{3}{$y=f(e^x+x^e)$。}
\end{minipage}\par\vfill
\noindent\begin{minipage}{\textwidth}\ProblemHeading{习题}{3.47}
设 $f(x)$ 处处可导，求下列各式的 $\dfrac{dy}{dx}$。
\SubQ{1}{$y=\sin(f(\sin(f(x))))$。}
\SubQ{2}{$y=f(f(f(\sin x+\cos x)))$。}
\SubQ{3}{$y=f(e^x)e^{f(x)}$。}
\end{minipage}\par\vfill
\noindent\begin{minipage}{\textwidth}\ProblemHeading{习题}{3.48}
求下列函数的导数，并指出不可导点。
\SubQ{1}{$\displaystyle y=\begin{cases}\dfrac{xe^{1/x}}{1+e^{1/x}},&x\ne0,\\0,&x=0.\end{cases}$}
\SubQ{2}{$y=|1-2x|\sin x$。}
\end{minipage}\par\vfill
\newpage
\noindent\begin{minipage}{\textwidth}\ProblemHeading{习题}{3.49}
设 $n$ 为正整数，$\displaystyle f(x)=\begin{cases}x^n\sin(1/x),&x\ne0,\\0,&x=0.\end{cases}$ 证明下列结论。
\SubQ{1}{当 $n=1$ 时，$f(x)$ 在 $x=0$ 处不可导。}
\SubQ{2}{当 $n=2$ 时，$f(x)$ 在 $x=0$ 处可导，但导函数在 $0$ 处不连续，且有第二类间断。}
\SubQ{3}{当 $n\ge3$ 时，$f(x)$ 在 $x=0$ 处可导，且导函数在 $0$ 处连续。}
\end{minipage}\par\vfill
\noindent\begin{minipage}{\textwidth}\ProblemHeading{习题}{3.50}
证明函数\[f(x)=\begin{cases}x^2\sin(1/x^2),&x\ne0,\\0,&x=0\end{cases}\]在 $[-1,1]$ 上处处可导，但其导函数在该区间上无界。
\end{minipage}\par\vfill
\noindent\begin{minipage}{\textwidth}\ProblemHeading{习题}{3.51（校订）}
在各函数导数非零的单调分支内，求反函数的微商 $dx/dy$，结果可用 $x$ 表示。
\SubQ{1}{$y=xe^x$（分别取 $x<-1$ 与 $x>-1$）。}
\SubQ{2}{$y=\arctan(1/x)$（$x\ne0$）。}
\SubQ{3}{$y=2x^{-x}-e^{-2x}$（$x>0$）。}
\end{minipage}\par\vfill
\noindent\begin{minipage}{\textwidth}\ProblemHeading{习题}{3.52（校订）}
在各函数导数非零的单调分支内，求反函数的微商 $dx/dy$，结果可用 $x$ 表示。
\SubQ{1}{$y=\ln(e^x+\sqrt{1+e^{2x}})$。}
\end{minipage}\par\vfill
\newpage
\noindent\begin{minipage}{\textwidth}\ProblemHeading{习题}{3.53}
证明可导偶函数的导数为奇函数，可导奇函数的导数为偶函数。
\end{minipage}\par\vfill
\noindent\begin{minipage}{\textwidth}\ProblemHeading{习题}{3.54}
证明可导周期函数的导数仍为周期函数。
\end{minipage}\par\vfill
\noindent\begin{minipage}{\textwidth}\ProblemHeading{习题}{3.55}
求下列各式的和，$n$ 为正整数。
\SubQ{1}{$P_n=1+2x+3x^2+\cdots+nx^{n-1}$。}
\SubQ{2}{$Q_n=1^2+2^2x+3^2x^2+\cdots+n^2x^{n-1}$。}
\SubQ{3}{$R_n=\cos1+2\cos2+\cdots+n\cos n$。}
\end{minipage}\par\vfill
\noindent\begin{minipage}{\textwidth}\ProblemHeading{习题}{3.56}
求下列函数的二阶导数，并指出二阶导数不存在的点。
\SubQ{1}{$y=e^{-x^2}$。}
\SubQ{2}{$y=x^2 2^x$。}
\SubQ{3}{$y=(1+x^2)\arctan x$。}
\end{minipage}\par\vfill
\newpage
\noindent\begin{minipage}{\textwidth}\ProblemHeading{习题}{3.57}
求下列函数的二阶导数，并指出二阶导数不存在的点。
\SubQ{1}{$y=\begin{cases}x^2,&x\ge0,\\-x^2,&x<0.\end{cases}$}
\end{minipage}\par\vfill
\noindent\begin{minipage}{\textwidth}\ProblemHeading{习题}{3.58}
设 $f(x)$ 处处有三阶导数，求下列函数的 $y''$ 和 $y'''$。
\SubQ{1}{$y=f(x^2)$。}
\SubQ{2}{$y=f(e^x+x)$。}
\end{minipage}\par\vfill
\noindent\begin{minipage}{\textwidth}\ProblemHeading{习题}{3.59}
设 $f(x)=x^n|x|$，其中 $n$ 为非负整数。证明 $f^{(n)}(0)$ 存在，而 $f^{(n+1)}(0)$ 不存在。
\end{minipage}\par\vfill
\noindent\begin{minipage}{\textwidth}\ProblemHeading{习题}{3.60}
求下列函数的 $n$ 阶导数，$n\ge1$。
\SubQ{1}{$x^2e^x$。}
\SubQ{2}{$(x^2+1)\sin x$。}
\SubQ{3}{$\displaystyle\frac1{x^2-3x+2}$。}
\end{minipage}\par\vfill
\newpage
\noindent\begin{minipage}{\textwidth}\ProblemHeading{习题}{3.61}
求下列函数的 $n$ 阶导数，$n\ge1$。
\SubQ{1}{$\sin x\cos x$。}
\end{minipage}\par\vfill
\noindent\begin{minipage}{\textwidth}\ProblemHeading{习题}{3.62}
求曲线 $y=\cos x$ 在 $x=\pi/4$ 处的切线方程。
\end{minipage}\par\vfill
\noindent\begin{minipage}{\textwidth}\ProblemHeading{习题}{3.63}
证明双曲线 $xy=1$ 上任一点处的切线与两坐标轴构成的三角形面积为定值。
\end{minipage}\par\vfill
\noindent\begin{minipage}{\textwidth}\ProblemHeading{习题}{3.64}
有一个底半径为 $r\,\mathrm{cm}$、高为 $h\,\mathrm{cm}$ 的正圆锥形容器，现以 $a\,\mathrm{cm}^3/\mathrm{s}$ 的速度自顶部向其内注水。求水面上升的速度。
\end{minipage}\par\vfill
\newpage
\noindent\begin{minipage}{\textwidth}\ProblemHeading{习题}{3.65}
水从高为 $18\,\mathrm{cm}$、底半径为 $6\,\mathrm{cm}$ 的圆锥形漏斗流入直径为 $10\,\mathrm{cm}$ 的圆柱形筒中。已知漏斗中水深为 $12\,\mathrm{cm}$ 时，水面下降速度为 $1\,\mathrm{cm}/\mathrm{min}$。求此时圆柱形筒中水面上升的速度。
\end{minipage}\par\vfill
\noindent\begin{minipage}{\textwidth}\ProblemHeading{例题}{3.66}
求 $y=e^{-ax}\sin bx$ 的微分，其中 $a,b$ 为常数。
\end{minipage}\par\vfill
\noindent\begin{minipage}{\textwidth}\ProblemHeading{例题}{3.67}
在定义域内求 $y=\ln(x+\sqrt{x^2+a^2})$ 的微分，其中 $a$ 为常数。
\end{minipage}\par\vfill
\noindent\begin{minipage}{\textwidth}\ProblemHeading{例题}{3.68}
设 $0<q<1$，函数 $y=y(x)$ 满足 $y-x-q\sin y=0$。求 $y$ 对 $x$ 的导数。
\end{minipage}\par\vfill
\newpage
\noindent\begin{minipage}{\textwidth}\ProblemHeading{习题}{3.69}
设 $y=x^2+x$。在 $x=1$ 处，对下列增量分别计算 $\Delta y$ 与 $dy$，并观察 $\Delta y-dy$ 随 $\Delta x$ 减小的变化。
\SubQ{1}{$\Delta x=10$。}
\SubQ{2}{$\Delta x=1$。}
\SubQ{3}{$\Delta x=0.1$。}
\end{minipage}\par\vfill
\noindent\begin{minipage}{\textwidth}\ProblemHeading{习题}{3.70}
设 $y=x^2+x$。在 $x=1$ 处，对下列增量分别计算 $\Delta y$ 与 $dy$，并观察 $\Delta y-dy$ 随 $\Delta x$ 减小的变化。
\SubQ{1}{$\Delta x=0.01$。}
\end{minipage}\par\vfill
\noindent\begin{minipage}{\textwidth}\ProblemHeading{习题}{3.71}
求下列函数的微分。
\SubQ{1}{$y=\ln(\pi/2-x/4)$。}
\SubQ{2}{$y=\sin x-x\cos x$。}
\SubQ{3}{$y=\arccos(1/|x|)$。}
\end{minipage}\par\vfill
\noindent\begin{minipage}{\textwidth}\ProblemHeading{习题}{3.72}
求下列函数的微分。
\SubQ{1}{$y=\ln\left|\dfrac{x-1}{x+1}\right|$。}
\SubQ{2}{$y=5^{\sqrt{\arctan(x^2)}}$。}
\SubQ{3}{$y=\tan^2(1+2x^2)$。}
\end{minipage}\par\vfill
\newpage
\noindent\begin{minipage}{\textwidth}\ProblemHeading{习题}{3.73}
求下列函数的微分。
\SubQ{1}{$y=e^{-x}\cos(3-x)$。}
\SubQ{2}{$y=x/\sqrt{x^2+1}$。}
\end{minipage}\par\vfill
\noindent\begin{minipage}{\textwidth}\ProblemHeading{习题}{3.74}
对下列参数方程，在 $dx/dt\ne0$（或 $dx/d\varphi\ne0$）处求 $dy/dx$ 与 $d^2y/dx^2$。
\SubQ{1}{$x=\ln(1+t^2),\ y=t-\arctan t$。}
\SubQ{2}{$x=t-\sin t,\ y=1-\cos t$。}
\SubQ{3}{$x=\varphi\cos\varphi,\ y=\varphi\sin\varphi$。}
\end{minipage}\par\vfill
\noindent\begin{minipage}{\textwidth}\ProblemHeading{习题}{3.75}
对下列参数方程，在 $dx/dt\ne0$（或 $dx/d\varphi\ne0$）处求 $dy/dx$ 与 $d^2y/dx^2$。
\SubQ{1}{$x=\cos^3\varphi,\ y=\sin^3\varphi$。}
\end{minipage}\par\vfill
\noindent\begin{minipage}{\textwidth}\ProblemHeading{习题}{3.76}
求下列参数曲线在指定参数值对应点处的切线方程。
\SubQ{1}{$x=\cos t,\ y=\sin t,\ t=\pi/4$。}
\SubQ{2}{$x=3t/(1+t^2),\ y=3t^2/(1+t^2),\ t=2$。}
\end{minipage}\par\vfill
\newpage
\noindent\begin{minipage}{\textwidth}\ProblemHeading{习题}{3.77}
Rolle 定理的条件为：函数在 $[a,b]$ 上连续、在 $(a,b)$ 内可导，且 $f(a)=f(b)$。分别举例说明，去掉下列任一条件，其结论“存在 $\xi\in(a,b)$ 使 $f^{\prime}(\xi)=0$”可能不成立。
\SubQ{1}{在闭区间上连续。}
\SubQ{2}{在开区间内可导。}
\SubQ{3}{两端点的函数值相等。}
\end{minipage}\par\vfill
\noindent\begin{minipage}{\textwidth}\ProblemHeading{例题}{3.78}
证明：对任意常数 $C$，方程 $x^3-3x+C=0$ 在 $[0,1]$ 中不可能有两个相异实根。
\end{minipage}\par\vfill
\noindent\begin{minipage}{\textwidth}\ProblemHeading{例题}{3.79}
设 $0<a<b$。证明\[\frac{b-a}{b}<\ln\frac ba<\frac{b-a}{a}.\]
\end{minipage}\par\vfill
\noindent\begin{minipage}{\textwidth}\ProblemHeading{例题}{3.80}
证明：当 $|x|\le1$ 时，$\arcsin x+\arccos x=\pi/2$。
\end{minipage}\par\vfill
\newpage
\noindent\begin{minipage}{\textwidth}\ProblemHeading{习题}{3.81}
设 $f(x)=(x-1)(x-2)(x-3)(x-4)$。确定 $f^{\prime}(x)=0$ 的实根个数及各根所在的区间。
\end{minipage}\par\vfill
\noindent\begin{minipage}{\textwidth}\ProblemHeading{习题}{3.82}
设 $f(x)$ 在 $[1,2]$ 上有二阶微商，且 $f(1)=f(2)=0$。令 $F(x)=(x-1)^2f(x)$，证明存在 $\xi\in(1,2)$ 使 $F^{\prime\prime}(\xi)=0$。
\end{minipage}\par\vfill
\noindent\begin{minipage}{\textwidth}\ProblemHeading{习题}{3.83}
举例说明下述中值定理的逆命题不成立：若 $\xi\in(a,b)$ 为指定一点，则存在不同的 $c,d\in[a,b]$，使\[\frac{f(c)-f(d)}{c-d}=f^{\prime}(\xi).\]
\end{minipage}\par\vfill
\noindent\begin{minipage}{\textwidth}\ProblemHeading{习题}{3.84}
证明下列不等式。
\SubQ{1}{若 $a>b>0,\ n>1$，则 $nb^{n-1}(a-b)<a^n-b^n<na^{n-1}(a-b)$。}
\SubQ{2}{若 $x>0$，则 $\dfrac{x}{1+x}<\ln(1+x)<x$。}
\SubQ{3}{若 $0<a<b$，则 $(a+b)\ln\dfrac{a+b}{2}<a\ln a+b\ln b$。}
\end{minipage}\par\vfill
\newpage
\noindent\begin{minipage}{\textwidth}\ProblemHeading{习题}{3.85}
证明下列不等式。
\SubQ{1}{若 $0<\alpha<\beta<\pi/2$，则 $\dfrac{\beta-\alpha}{\cos^2\alpha}<\tan\beta-\tan\alpha<\dfrac{\beta-\alpha}{\cos^2\beta}$。}
\end{minipage}\par\vfill
\noindent\begin{minipage}{\textwidth}\ProblemHeading{习题}{3.86}
证明下列恒等式。
\SubQ{1}{$\arctan x=\arcsin\dfrac{x}{\sqrt{1+x^2}}\quad(x\in\mathbb R)$。}
\SubQ{2}{$\arctan x+\arctan\dfrac{1-x}{1+x}=\begin{cases}\pi/4,&x>-1,\\-3\pi/4,&x<-1.\end{cases}$}
\end{minipage}\par\vfill
\noindent\begin{minipage}{\textwidth}\ProblemHeading{习题}{3.87}
设 $f(x)$ 在 $[0,1]$ 上可微，对每个 $x\in[0,1]$ 都有 $f(x)\in(0,1)$ 且 $f^{\prime}(x)\ne1$。证明存在唯一的 $\xi\in(0,1)$ 使 $f(\xi)=\xi$。
\end{minipage}\par\vfill
\noindent\begin{minipage}{\textwidth}\ProblemHeading{习题}{3.88}
设 $f(x)$ 在 $[0,1]$ 上连续，在 $(0,1)$ 内可微且 $|f^{\prime}(x)|<1$，又有 $f(0)=f(1)$。证明对任意 $x_1,x_2\in[0,1]$，都有 $|f(x_1)-f(x_2)|<1/2$。
\end{minipage}\par\vfill
\newpage
\noindent\begin{minipage}{\textwidth}\ProblemHeading{习题}{3.89}
设 $f(x)$ 在实数轴上处处可导，并满足 $f^{\prime}(x)=f(x)$。证明 $f(x)=Ce^x$，其中 $C$ 为常数。
\end{minipage}\par\vfill
\noindent\begin{minipage}{\textwidth}\ProblemHeading{习题}{3.90}
设非常数函数 $f(x)$ 在 $[a,b]$ 上连续，在 $(a,b)$ 内可微，且 $f(a)=f(b)$。证明存在 $\xi\in(a,b)$ 使 $f^{\prime}(\xi)>0$。
\end{minipage}\par\vfill
\noindent\begin{minipage}{\textwidth}\ProblemHeading{习题}{3.91}
设 $f(x)$ 在 $[a,+\infty)$ 上可微，且 $\lim_{x\to+\infty}f^{\prime}(x)=0$。证明：
\SubQ{1}{$\displaystyle\lim_{x\to+\infty}\bigl(f(x+1)-f(x)\bigr)=0$。}
\SubQ{2}{$\displaystyle\lim_{x\to+\infty}\frac{f(x)}x=0$。}
\end{minipage}\par\vfill
\noindent\begin{minipage}{\textwidth}\ProblemHeading{习题}{3.92}
考察可导函数及其导函数的有界性。
\SubQ{1}{证明：若 $f(x)$ 在有限开区间 $(a,b)$ 内的导数有界，则 $f(x)$ 有界。}
\SubQ{2}{当区间为无限区间时，上述结论是否仍成立？}
\SubQ{3}{上述有限区间情形的逆命题是否成立？}
\end{minipage}\par\vfill
\newpage
\noindent\begin{minipage}{\textwidth}\ProblemHeading{习题}{3.93}
设存在常数 $M$，使任意实数 $x,y$ 都满足 $|f(x)-f(y)|\le M|x-y|^2$。证明 $f(x)$ 为常数。
\end{minipage}\par\vfill
\noindent\begin{minipage}{\textwidth}\ProblemHeading{习题}{3.94}
设函数 $f(x)$ 在区间 $I$ 上连续，在 $I$ 内除至多有限个点外处处有导数。证明：若这些点处的导数均为正（均为负），则 $f(x)$ 在 $I$ 上严格递增（严格递减）。
\end{minipage}\par\vfill
\noindent\begin{minipage}{\textwidth}\ProblemHeading{习题}{3.95}
设 $f(x),g(x)$ 在区间 $I$ 上连续，在 $I$ 内除至多有限个点外处处可导且 $f^{\prime}(x)>g^{\prime}(x)$。设 $a$ 是 $I$ 的一个非端点，且 $f(a)=g(a)$。证明：当 $x>a$ 时 $f(x)>g(x)$；当 $x<a$ 时 $f(x)<g(x)$。
\end{minipage}\par\vfill
\noindent\begin{minipage}{\textwidth}\ProblemHeading{习题}{3.96}
设 $f(x)$ 在 $[0,+\infty)$ 上可微，$f(0)=0$，且 $f^{\prime}(x)$ 严格递增。证明 $f(x)/x$ 在 $(0,+\infty)$ 内严格递增。
\end{minipage}\par\vfill
\newpage
\noindent\begin{minipage}{\textwidth}\ProblemHeading{习题}{3.97}
设 $x_0$ 为函数 $f$ 的驻点，即 $f^{\prime}(x_0)=0$，且 $f$ 在 $x_0$ 处有二阶导数。完成以下各问。
\SubQ{1}{若 $f^{\prime\prime}(x_0)<0$，证明 $x_0$ 为局部极大值点。}
\SubQ{2}{若 $f^{\prime\prime}(x_0)>0$，证明 $x_0$ 为局部极小值点。}
\SubQ{3}{举例说明：当 $f^{\prime\prime}(x_0)=0$ 时，$x_0$ 可以是局部极大值点。}
\end{minipage}\par\vfill
\noindent\begin{minipage}{\textwidth}\ProblemHeading{习题}{3.98}
设 $x_0$ 为函数 $f$ 的驻点，即 $f^{\prime}(x_0)=0$，且 $f$ 在 $x_0$ 处有二阶导数。完成以下各问。
\SubQ{1}{举例说明：当 $f^{\prime\prime}(x_0)=0$ 时，$x_0$ 可以是局部极小值点。}
\SubQ{2}{举例说明：当 $f^{\prime\prime}(x_0)=0$ 时，$x_0$ 可以不是极值点。}
\end{minipage}\par\vfill
\noindent\begin{minipage}{\textwidth}\ProblemHeading{习题}{3.99}
设 $f(x)$ 在 $[0,1]$ 上有二阶导数，且 $f(0)=f^{\prime}(0)$、$f(1)=f^{\prime}(1)$。证明存在 $\xi\in(0,1)$ 使 $f(\xi)=f^{\prime\prime}(\xi)$。
\end{minipage}\par\vfill
\noindent\begin{minipage}{\textwidth}\ProblemHeading{习题}{3.100}
求下列函数的单调区间与极值。
\SubQ{1}{$2x^3-3x^2$。}
\SubQ{2}{$x^{2/3}$。}
\SubQ{3}{$x^2e^{-x^2}$。}
\end{minipage}\par\vfill
\newpage
\noindent\begin{minipage}{\textwidth}\ProblemHeading{习题}{3.101}
求下列函数的单调区间与极值。
\SubQ{1}{$x^{1/x}\ (x>0)$。}
\SubQ{2}{$\dfrac{(\ln x)^2}{x}\ (x>0)$。}
\SubQ{3}{$\arctan x-\dfrac12\ln(1+x^2)$。}
\end{minipage}\par\vfill
\noindent\begin{minipage}{\textwidth}\ProblemHeading{习题}{3.102}
求下列函数在指定区间上的最大值和最小值；不存在时说明理由。
\SubQ{1}{$x^4-2x^2+5,\quad[-2,2]$。}
\SubQ{2}{$\sin2x-x,\quad[-\pi/2,\pi/2]$。}
\SubQ{3}{$\arctan\dfrac{1-x}{1+x},\quad[0,1]$。}
\end{minipage}\par\vfill
\noindent\begin{minipage}{\textwidth}\ProblemHeading{习题}{3.103}
求下列函数在指定区间上的最大值和最小值；不存在时说明理由。
\SubQ{1}{$x\ln x,\quad(0,+\infty)$。}
\end{minipage}\par\vfill
\noindent\begin{minipage}{\textwidth}\ProblemHeading{习题}{3.104}
证明下列不等式。
\SubQ{1}{$\dfrac1{2^{p-1}}\le x^p+(1-x)^p\le1,\quad x\in(0,1],\ p>1$。}
\SubQ{2}{$\tan x>x-\dfrac{x^3}{3},\quad0<x<\pi/2$。}
\SubQ{3}{$\dfrac{\tan x_2}{\tan x_1}>\dfrac{x_2}{x_1},\quad0<x_1<x_2<\pi/2$。}
\end{minipage}\par\vfill
\newpage
\noindent\begin{minipage}{\textwidth}\ProblemHeading{习题}{3.105}
证明下列不等式。
\SubQ{1}{$\ln(1+x)>\dfrac{\arctan x}{1+x},\quad x>0$。}
\SubQ{2}{$1+x\ln(x+\sqrt{1+x^2})\ge\sqrt{1+x^2},\quad x\in\mathbb R$。}
\SubQ{3}{$\dfrac{x}{\sin x}>\dfrac43-\dfrac13\cos x,\quad0<x<\pi/2$；并证明常数 $4/3$ 不能增大。}
\end{minipage}\par\vfill
\noindent\begin{minipage}{\textwidth}\ProblemHeading{习题}{3.106}
证明下列不等式。
\SubQ{1}{$\left(1-\dfrac1x\right)^{x-1}\left(1+\dfrac1x\right)^{x+1}<4,\quad x>1$。}
\SubQ{2}{$x^{a-1}+x^{a+1}\ge\left(\dfrac{1-a}{1+a}\right)^{(a-1)/2}+\left(\dfrac{1-a}{1+a}\right)^{(a+1)/2},\quad0<x<1,\ 0<a<1$。}
\end{minipage}\par\vfill
\noindent\begin{minipage}{\textwidth}\ProblemHeading{习题}{3.107}
确定下列函数的实零点个数及所在范围。
\SubQ{1}{$x^3-6x^2+9x-10$。}
\SubQ{2}{$ax-\ln x$，其中 $a>0$、$x>0$。}
\end{minipage}\par\vfill
\noindent\begin{minipage}{\textwidth}\ProblemHeading{习题}{3.108}
设 $a\in(0,1)$，$b_1=1-a$，且\[b_{n+1}=\frac{b_n}{1-e^{-b_n}}-a,\qquad n=1,2,\ldots.\]判断数列 $\{b_n\}$ 是否收敛；若不收敛，给出证明；若收敛，求其极限。
\end{minipage}\par\vfill
\newpage
\noindent\begin{minipage}{\textwidth}\ProblemHeading{例题}{3.109}
求 $\displaystyle\lim_{x\to0}\frac{e^{2x}-1}{\ln(1+x)}$。
\end{minipage}\par\vfill
\noindent\begin{minipage}{\textwidth}\ProblemHeading{例题}{3.110}
求 $\displaystyle\lim_{x\to0}\frac{e^x+e^{-x}-2}{\sin^2x}$。
\end{minipage}\par\vfill
\noindent\begin{minipage}{\textwidth}\ProblemHeading{例题}{3.111}
设 $\alpha>0$。求 $\displaystyle\lim_{x\to+\infty}\frac{\ln x}{x^\alpha}$。
\end{minipage}\par\vfill
\noindent\begin{minipage}{\textwidth}\ProblemHeading{例题}{3.112}
设 $\mu>0,\ a>1$。求 $\displaystyle\lim_{x\to+\infty}\frac{x^\mu}{a^x}$。
\end{minipage}\par\vfill
\newpage
\noindent\begin{minipage}{\textwidth}\ProblemHeading{例题}{3.113}
求 $\displaystyle\lim_{x\to+\infty}x\left(\frac\pi2-\arctan x\right)$。
\end{minipage}\par\vfill
\noindent\begin{minipage}{\textwidth}\ProblemHeading{例题}{3.114}
求 $\displaystyle\lim_{x\to1}\left(\frac1{\ln x}+\frac1{1-x}\right)$。
\end{minipage}\par\vfill
\noindent\begin{minipage}{\textwidth}\ProblemHeading{例题}{3.115}
求 $\displaystyle\lim_{x\to0}\left(\frac{\sin x}{x}\right)^{1/x^2}$。
\end{minipage}\par\vfill
\noindent\begin{minipage}{\textwidth}\ProblemHeading{例题}{3.116}
求 $\displaystyle\lim_{x\to0^+}x^x$。
\end{minipage}\par\vfill
\newpage
\noindent\begin{minipage}{\textwidth}\ProblemHeading{例题}{3.117}
求 $\displaystyle\lim_{x\to+\infty}x^{1/\sqrt{x}}$。
\end{minipage}\par\vfill
\noindent\begin{minipage}{\textwidth}\ProblemHeading{例题}{3.118}
求 $\displaystyle\lim_{x\to0}(1-\cos x)^{\tan x}$。
\end{minipage}\par\vfill
\noindent\begin{minipage}{\textwidth}\ProblemHeading{例题}{3.119}
求 $\displaystyle\lim_{x\to+\infty}x^{\arctan x-\pi/2}$。
\end{minipage}\par\vfill
\noindent\begin{minipage}{\textwidth}\ProblemHeading{例题}{3.120}
求 $\displaystyle\lim_{x\to0}\left(\frac1{x^2}-\frac1{x\tan x}\right)$。
\end{minipage}\par\vfill
\newpage
\noindent\begin{minipage}{\textwidth}\ProblemHeading{习题}{3.121}
Cauchy 中值定理的一种形式为：若 $f,g$ 在 $[a,b]$ 上连续、在 $(a,b)$ 内可导，且 $g^{\prime}(x)\ne0$，则存在 $\xi\in(a,b)$ 使\[\frac{f(b)-f(a)}{g(b)-g(a)}=\frac{f^{\prime}(\xi)}{g^{\prime}(\xi)}.\]给出该定理的几何解释。
\end{minipage}\par\vfill
\noindent\begin{minipage}{\textwidth}\ProblemHeading{习题}{3.122（校订）}
设 $f(x)=x^2$、$g(x)=x^3$。说明为何不能在 $[-1,1]$ 上直接套用 Cauchy 中值定理的商形式\[\frac{f(1)-f(-1)}{g(1)-g(-1)}=\frac{f^{\prime}(\xi)}{g^{\prime}(\xi)},\qquad -1<\xi<1.\]
\end{minipage}\par\vfill
\noindent\begin{minipage}{\textwidth}\ProblemHeading{习题}{3.123}
设 $b>a>0$，函数 $f$ 在 $[a,b]$ 上连续，在 $(a,b)$ 内可微。证明存在 $\xi\in(a,b)$，使 $2\xi(f(b)-f(a))=(b^2-a^2)f^{\prime}(\xi)$。
\end{minipage}\par\vfill
\noindent\begin{minipage}{\textwidth}\ProblemHeading{习题}{3.124}
设 $f$ 在 $[a,b]$ 上连续，其中 $ab>0$，并在 $(a,b)$ 内可微。证明存在 $\xi\in(a,b)$，使\[\frac{af(b)-bf(a)}{a-b}=f(\xi)-\xi f^{\prime}(\xi).\]
\end{minipage}\par\vfill
\newpage
\noindent\begin{minipage}{\textwidth}\ProblemHeading{习题}{3.125}
求下列极限。
\SubQ{1}{$\displaystyle\lim_{x\to0}\frac{\sqrt[m]{1+\alpha x}-\sqrt[n]{1+\beta x}}x$，其中 $m,n$ 为正整数，$\alpha,\beta$ 为实数。}
\SubQ{2}{$\displaystyle\lim_{x\to0}\frac{(1+mx)^n-(1+nx)^m}{x^2}$，其中 $m,n$ 为自然数。}
\SubQ{3}{$\displaystyle\lim_{x\to1}\frac{x^3+x-2}{x^2-3x+2}$。}
\end{minipage}\par\vfill
\noindent\begin{minipage}{\textwidth}\ProblemHeading{习题}{3.126}
求下列极限。
\SubQ{1}{$\displaystyle\lim_{x\to0}\frac{x-\arcsin x}{\sin^3x}$。}
\SubQ{2}{$\displaystyle\lim_{x\to0}\frac{e^x-1}{\sin x}$。}
\SubQ{3}{$\displaystyle\lim_{x\to0}\frac{(1+x)^\alpha-1}x$，其中 $\alpha\in\mathbb R$。}
\end{minipage}\par\vfill
\noindent\begin{minipage}{\textwidth}\ProblemHeading{习题}{3.127}
求下列极限。
\SubQ{1}{$\displaystyle\lim_{x\to0}\frac{e^{-1/x^2}}x$。}
\SubQ{2}{$\displaystyle\lim_{x\to0}\frac{(a+x)^x-a^x}{x^2}$，其中 $a>0$。}
\SubQ{3}{$\displaystyle\lim_{x\to0}x\sin\frac1x$。}
\end{minipage}\par\vfill
\noindent\begin{minipage}{\textwidth}\ProblemHeading{习题}{3.128}
求下列极限。
\SubQ{1}{$\displaystyle\lim_{x\to0}\frac{x^2\cos(1/x)}{\tan x}$。}
\SubQ{2}{$\displaystyle\lim_{x\to0}\left(\frac1{\arctan^2x}-\frac1{x^2}\right)$。}
\SubQ{3}{$\displaystyle\lim_{x\to1^-}\ln x\ln(1-x)$。}
\end{minipage}\par\vfill
\newpage
\noindent\begin{minipage}{\textwidth}\ProblemHeading{习题}{3.129}
求下列极限。
\SubQ{1}{$\displaystyle\lim_{x\to(\pi/2)^-}(\tan x)^{2x-\pi}$。}
\SubQ{2}{$\displaystyle\lim_{x\to0}\left(\frac{(1+x)^{1/x}}e\right)^{1/x}$。}
\SubQ{3}{$\displaystyle\lim_{x\to1^-}\frac{\ln(1-x)+\tan(\pi x/2)}{\cot(\pi x)}$。}
\end{minipage}\par\vfill
\noindent\begin{minipage}{\textwidth}\ProblemHeading{习题}{3.130}
求下列极限。
\SubQ{1}{$\displaystyle\lim_{x\to0}\frac{(\arcsin(x^3))^2}{(1-\cos x)(e^{x^2}-1)\tan^2x}$。}
\SubQ{2}{$\displaystyle\lim_{x\to+\infty}\left(\frac{\ln(1+x)}x\right)^{1/x}$。}
\SubQ{3}{$\displaystyle\lim_{x\to0}\frac{\arctan(x^2)}{\sqrt{1+x\sin x}-\sqrt{\cos x}}\left(2-\frac{x}{e^x-1}\right)$。}
\end{minipage}\par\vfill
\noindent\begin{minipage}{\textwidth}\ProblemHeading{习题}{3.131}
求下列极限。
\SubQ{1}{$\displaystyle\lim_{n\to\infty}\frac{n^k}{a^n}$，其中 $a>1,\ k>0$。}
\SubQ{2}{$\displaystyle\lim_{n\to\infty}\frac{\ln n}{n^k}$，其中 $k>0$。}
\end{minipage}\par\vfill
\noindent\begin{minipage}{\textwidth}\ProblemHeading{习题}{3.132}
设 $f(x)$ 在 $[0,a]$ 上有二阶连续导数，$f^{\prime}(0)=1$、$f^{\prime\prime}(0)\ne0$，且当 $0<x<a$ 时 $0<f(x)<x$。令 $x_1\in(0,a)$、$x_{n+1}=f(x_n)$。
\SubQ{1}{证明 $\{x_n\}$ 收敛，并求其极限。}
\SubQ{2}{判断 $\{nx_n\}$ 是否收敛；若收敛，求其极限。}
\end{minipage}\par\vfill
\newpage
\noindent\begin{minipage}{\textwidth}\ProblemHeading{例题}{3.133}
证明：当 $0<x\le\pi/2$ 时，$\dfrac2\pi x\le\sin x$。
\end{minipage}\par\vfill
\noindent\begin{minipage}{\textwidth}\ProblemHeading{例题}{3.134}
求函数 $f(x)=e^{-x^2}$ 在实数轴上的单调区间与极值。
\end{minipage}\par\vfill
\noindent\begin{minipage}{\textwidth}\ProblemHeading{例题}{3.135}
求函数 $f(x)=x^3-3x^2-9x+5$ 在区间 $[-4,4]$ 上的最大值与最小值。
\end{minipage}\par\vfill
\noindent\begin{minipage}{\textwidth}\ProblemHeading{例题}{3.136（校订）}
设 $x>0,\ 0<\alpha<1$。证明 $x^\alpha-\alpha x\le1-\alpha$。
\end{minipage}\par\vfill
\newpage
\noindent\begin{minipage}{\textwidth}\ProblemHeading{例题}{3.137}
在给定直线上求一点 $P$，使它到直线外两个已知点 $A,B$ 的距离之和最小。
\end{minipage}\par\vfill
\noindent\begin{minipage}{\textwidth}\ProblemHeading{例题}{3.138}
证明函数 $f(x)=x^2$ 在 $(-\infty,+\infty)$ 上为凸函数。
\end{minipage}\par\vfill
\noindent\begin{minipage}{\textwidth}\ProblemHeading{习题}{3.139}
设 $f$ 为区间 $I$ 上的连续函数，在 $I$ 的内部二阶可导。证明以下判别条件。
\SubQ{1}{$f$ 在 $I$ 上为凸函数，当且仅当在 $I$ 内部处处有 $f^{\prime\prime}(x)\ge0$。}
\SubQ{2}{$f$ 在 $I$ 上严格凸，当且仅当 $f^{\prime\prime}(x)\ge0$，且在任一非退化子区间上 $f^{\prime\prime}$ 不恒为零。}
\end{minipage}\par\vfill
\noindent\begin{minipage}{\textwidth}\ProblemHeading{例题}{3.140}
确定函数 $f(x)=e^{-x^2}$ 的图像在 $(-\infty,+\infty)$ 上的形态。
\end{minipage}\par\vfill
\newpage
\noindent\begin{minipage}{\textwidth}\ProblemHeading{例题}{3.141}
计算抛物线 $y=\dfrac12ax^2$ 的有向曲率，其中 $a\ne0$，取 $x$ 增大的方向。
\end{minipage}\par\vfill
\noindent\begin{minipage}{\textwidth}\ProblemHeading{例题}{3.142}
计算椭圆 $x=a\cos t,\ y=b\sin t,\ 0\le t\le2\pi$ 上各点的有向曲率，其中 $a,b>0$，取 $t$ 增大的方向。
\end{minipage}\par\vfill
\noindent\begin{minipage}{\textwidth}\ProblemHeading{习题}{3.143}
证明 Jensen 不等式：若 $f(x)$ 为区间 $I$ 上的凸函数，$x_1,\ldots,x_n\in I$，且 $\alpha_1,\ldots,\alpha_n>0$、$\sum_{i=1}^n\alpha_i=1$，则\[f\left(\sum_{i=1}^n\alpha_i x_i\right)\le\sum_{i=1}^n\alpha_i f(x_i).\]
\end{minipage}\par\vfill
\noindent\begin{minipage}{\textwidth}\ProblemHeading{习题}{3.144}
设 $x_1,\ldots,x_n$ 和 $\lambda_1,\ldots,\lambda_n$ 都为正数，且 $\sum_{i=1}^n\lambda_i=1$。证明加权算术—几何平均不等式\[x_1^{\lambda_1}\cdots x_n^{\lambda_n}\le\lambda_1x_1+\cdots+\lambda_nx_n,\]并由此推出 $\sqrt[n]{x_1\cdots x_n}\le(x_1+\cdots+x_n)/n$。
\end{minipage}\par\vfill
\newpage
\noindent\begin{minipage}{\textwidth}\ProblemHeading{习题}{3.145}
设 $a,b,c,d$ 满足 $a/b\le c/d$，其中 $b>0,d>0$。证明\[\frac ab\le\frac{a+c}{b+d}\le\frac cd.\]
\end{minipage}\par\vfill
\noindent\begin{minipage}{\textwidth}\ProblemHeading{习题}{3.146}
设 $f(x)$ 为区间 $I$ 上的凸函数。证明 $f$ 在 $I$ 的每个内点连续。
\end{minipage}\par\vfill
\noindent\begin{minipage}{\textwidth}\ProblemHeading{习题}{3.147}
设函数 $f(x)$ 在区间 $I$ 上处处可导，除有限个点外均有 $f^{\prime\prime}(x)>0$。证明 $f$ 在 $I$ 上为凸函数，并将结论应用于 $-f$，得到凹函数的类似判别。
\end{minipage}\par\vfill
\noindent\begin{minipage}{\textwidth}\ProblemHeading{习题}{3.148}
设 $f$ 在区间 $I$ 内有连续的二阶导数，$x_0$ 为 $f$ 的一个拐点。证明 $f^{\prime\prime}(x_0)=0$。
\end{minipage}\par\vfill
\newpage
\noindent\begin{minipage}{\textwidth}\ProblemHeading{习题}{3.149}
设 $f$ 在 $x_0$ 及其附近二阶可导，$f^{\prime\prime}(x_0)=0$，且 $f^{(3)}(x_0)$ 存在且不为零。证明 $x_0$ 为 $f$ 的拐点。
\end{minipage}\par\vfill
\noindent\begin{minipage}{\textwidth}\ProblemHeading{习题}{3.150}
求下列函数的凸区间、凹区间和拐点。
\SubQ{1}{$y=2x^3-3x^2-36x+25$。}
\SubQ{2}{$y=x+1/x$。}
\SubQ{3}{$y=x^{5/3}$。}
\end{minipage}\par\vfill
\noindent\begin{minipage}{\textwidth}\ProblemHeading{习题}{3.151}
求下列函数的凸区间、凹区间和拐点。
\SubQ{1}{$y=(1+x^2)e^x$。}
\SubQ{2}{$y=x^4$。}
\SubQ{3}{$y=x+\sin x$。}
\end{minipage}\par\vfill
\noindent\begin{minipage}{\textwidth}\ProblemHeading{习题}{3.152}
求 $a,b$ 的值，使点 $(1,3)$ 为曲线 $y=ax^3+bx^2$ 的拐点。
\end{minipage}\par\vfill
\newpage
\noindent\begin{minipage}{\textwidth}\ProblemHeading{习题}{3.153}
描绘下列各曲线的图形。
\SubQ{1}{$y=x^3+6x^2-15x-20$。}
\SubQ{2}{$y=\dfrac{x^3}{2(1+x)^2}$。}
\SubQ{3}{$y=x-2\arctan x$。}
\end{minipage}\par\vfill
\noindent\begin{minipage}{\textwidth}\ProblemHeading{习题}{3.154}
描绘下列各曲线的图形。
\SubQ{1}{$y=xe^{-x}$。}
\end{minipage}\par\vfill
\noindent\begin{minipage}{\textwidth}\ProblemHeading{习题}{3.155（校订）}
曲线在点 $M$ 处的有向曲率为 $\kappa\ne0$，曲率半径定义为 $\rho=1/|\kappa|$；曲率中心在该点法线上、曲线弯曲中心的一侧，与 $M$ 的距离为 $\rho$。求下列曲线在指定点的有向曲率、曲率中心与曲率半径，方向取 $x$ 增大的方向。
\SubQ{1}{$xy=1$，在点 $(1,1)$ 处。}
\SubQ{2}{$y=e^{-x^2}$，在点 $(0,1)$ 处。}
\end{minipage}\par\vfill
\noindent\begin{minipage}{\textwidth}\ProblemHeading{习题}{3.156}
求下列参数曲线在指定参数值处的有向曲率，方向取 $t$ 增大的方向。
\SubQ{1}{$x=3t^2,\ y=3t-t^3,\quad t=1$。}
\SubQ{2}{$x=\cos t+t\sin t,\ y=\sin t-t\cos t,\quad t=\pi/2$。}
\end{minipage}\par\vfill
\newpage
\noindent\begin{minipage}{\textwidth}\ProblemHeading{习题}{3.157}
对数曲线 $y=\ln x$ 上哪一点的曲率半径最小？求出该点的曲率半径。
\end{minipage}\par\vfill
\noindent\begin{minipage}{\textwidth}\ProblemHeading{习题}{3.158}
设 $f(x)$ 为定义在 $(-\infty,+\infty)$ 上的凸函数，且有上界。证明 $f$ 为常数。
\end{minipage}\par\vfill
\noindent\begin{minipage}{\textwidth}\ProblemHeading{例题}{3.159}
求函数 $e^{-x^4}$ 的 $4n$ 次 Maclaurin 公式，其中 $n$ 为正整数。
\end{minipage}\par\vfill
\noindent\begin{minipage}{\textwidth}\ProblemHeading{例题}{3.160}
求 $\cos^2x$ 的 $2n$ 次 Maclaurin 公式，其中 $n$ 为正整数。
\end{minipage}\par\vfill
\newpage
\noindent\begin{minipage}{\textwidth}\ProblemHeading{例题}{3.161}
求 $\displaystyle\lim_{x\to0}\frac{e^{-x^4}-\cos^2x-x^2}{x^4}$。
\end{minipage}\par\vfill
\noindent\begin{minipage}{\textwidth}\ProblemHeading{例题}{3.162}
设 $f(x)=1/x$。求 $f$ 在 $x=2$ 处的 Taylor 公式。
\end{minipage}\par\vfill
\noindent\begin{minipage}{\textwidth}\ProblemHeading{例题}{3.163}
计算 $e$ 的近似值，使误差不超过 $10^{-5}$。
\end{minipage}\par\vfill
\noindent\begin{minipage}{\textwidth}\ProblemHeading{例题}{3.164}
证明 $e$ 为无理数。
\end{minipage}\par\vfill
\newpage
\noindent\begin{minipage}{\textwidth}\ProblemHeading{例题}{3.165}
设 $f$ 在 $[0,1]$ 上二阶可导，任意 $x\in[0,1]$ 都满足 $|f^{\prime\prime}(x)|\le M$，其中 $M$ 为常数，且 $f(0)=f(1)=0$。证明在 $[0,1]$ 上有 $|f(x)|\le M/8$。
\end{minipage}\par\vfill
\noindent\begin{minipage}{\textwidth}\ProblemHeading{习题}{3.166}
写出下列函数具有 Peano 余项的 Maclaurin 展开式。
\SubQ{1}{$y=\dfrac{x^3+2x+1}{x-1}$。}
\SubQ{2}{$y=\sin^2x$。}
\end{minipage}\par\vfill
\noindent\begin{minipage}{\textwidth}\ProblemHeading{习题}{3.167}
求 $e^{\sin x}$ 具有 Peano 余项的三阶 Maclaurin 展开式。
\end{minipage}\par\vfill
\noindent\begin{minipage}{\textwidth}\ProblemHeading{习题}{3.168}
求 $\ln(\cos x)$ 具有 Peano 余项的六阶 Maclaurin 展开式。
\end{minipage}\par\vfill
\newpage
\noindent\begin{minipage}{\textwidth}\ProblemHeading{习题}{3.169}
已知 $f(x)$ 为四次多项式，且 $f(2)=-1$、$f^{\prime}(2)=0$、$f^{\prime\prime}(2)=2$、$f^{(3)}(2)=-12$、$f^{(4)}(2)=24$。计算 $f(-1)$、$f^{\prime}(0)$、$f^{\prime\prime}(1)$。
\end{minipage}\par\vfill
\noindent\begin{minipage}{\textwidth}\ProblemHeading{习题}{3.170}
求下列具有 Lagrange 余项的 Taylor 展开式。
\SubQ{1}{$y=\tan x$ 在 $x=0$ 处的二阶展开式。}
\SubQ{2}{$y=1/x$ 在 $x=-1$ 处的 $n$ 阶展开式。}
\end{minipage}\par\vfill
\noindent\begin{minipage}{\textwidth}\ProblemHeading{习题}{3.171}
求下列极限。
\SubQ{1}{$\displaystyle\lim_{x\to0}\frac{\cos x-e^{-x^2/2}}{\sin^4x}$。}
\SubQ{2}{$\displaystyle\lim_{x\to0}\frac{\sqrt[4]{1+x^2}-\sqrt[4]{1-x^2}}{x^2}$。}
\SubQ{3}{$\displaystyle\lim_{x\to\infty}\left[x-x^2\ln\left(1+\frac1x\right)\right]$。}
\end{minipage}\par\vfill
\noindent\begin{minipage}{\textwidth}\ProblemHeading{习题}{3.172}
求下列极限。
\SubQ{1}{$\displaystyle\lim_{x\to0}\frac{\cos(\sin x)-\cos x}{\sin^4x}$。}
\end{minipage}\par\vfill
\newpage
\noindent\begin{minipage}{\textwidth}\ProblemHeading{习题}{3.173}
设 $f$ 处处有 $n+1$ 阶导数。证明：$f(x)$ 为次数不超过 $n$ 的多项式，当且仅当 $f^{(n+1)}(x)$ 恒为零。
\end{minipage}\par\vfill
\noindent\begin{minipage}{\textwidth}\ProblemHeading{习题}{3.174}
设 $f$ 在 $[0,2]$ 上二阶可导，且对任意 $x\in[0,2]$ 有 $|f(x)|\le1$、$|f^{\prime\prime}(x)|\le1$。证明 $|f^{\prime}(x)|\le2$ 对所有 $x\in[0,2]$ 成立。
\end{minipage}\par\vfill
\noindent\begin{minipage}{\textwidth}\ProblemHeading{习题}{3.175}
设 $n$ 为正整数，定义\[f(x)=\begin{cases}x^{n+1},&x\in\mathbb Q,\\0,&x\notin\mathbb Q.\end{cases}\]证明 $f^{\prime}(0)=0$，但 $f^{\prime\prime}(0)$ 不存在。
\end{minipage}\par\vfill
\noindent\begin{minipage}{\textwidth}\ProblemHeading{习题}{3.176}
定义\[f(x)=\begin{cases}e^{-1/x^2},&x\ne0,\\0,&x=0.\end{cases}\]证明 $f$ 在 $x=0$ 处的任意阶导数都存在，且都等于零。
\end{minipage}\par\vfill
\newpage
\noindent\begin{minipage}{\textwidth}\ProblemHeading{习题}{3.177}
设函数 $f$ 在驻点 $x_0$ 处的 $n$ 阶微商存在，$n\ge2$，并且 $f^{\prime}(x_0)=\cdots=f^{(n-1)}(x_0)=0$、$f^{(n)}(x_0)\ne0$。证明：
\SubQ{1}{若 $n$ 为奇数，则 $f$ 在 $x_0$ 处无极值。}
\SubQ{2}{若 $n$ 为偶数，则 $f$ 在 $x_0$ 处取极值；当 $f^{(n)}(x_0)>0$ 时取极小值，当 $f^{(n)}(x_0)<0$ 时取极大值。}
\end{minipage}\par\vfill
\noindent\begin{minipage}{\textwidth}\ProblemHeading{习题}{3.178}
设 $f(x)=x(x+1)(x+2)\cdots(x+n)$，其中 $n$ 为正整数。求 $f^{\prime}(0)$。
\end{minipage}\par\vfill
\noindent\begin{minipage}{\textwidth}\ProblemHeading{习题}{3.179}
设奇函数 $f$ 在 $(-\infty,+\infty)$ 上有二阶连续导数，定义\[g(x)=\begin{cases}f(x)/x,&x\ne0,\\a,&x=0.\end{cases}\]
\SubQ{1}{确定 $a$ 的值，使 $g$ 在整个实轴上连续。}
\SubQ{2}{对使 $g$ 连续的 $a$，证明 $g$ 在整个实轴上可导，且导函数连续。}
\end{minipage}\par\vfill
\noindent\begin{minipage}{\textwidth}\ProblemHeading{习题}{3.180}
设实数 $a_0,\ldots,a_n$ 满足 $\dfrac{a_0}{n+1}+\dfrac{a_1}{n}+\cdots+a_n=0$。证明方程 $a_0x^n+a_1x^{n-1}+\cdots+a_n=0$ 在 $(0,1)$ 内至少有一个根。
\end{minipage}\par\vfill
\newpage
\noindent\begin{minipage}{\textwidth}\ProblemHeading{习题}{3.181}
设 $f$ 在 $[a,b]$ 上连续，在 $(a,b)$ 内可微，且 $f(a)=f(b)=0$。证明存在 $\xi\in(a,b)$，使 $f^{\prime}(\xi)+f(\xi)=0$。
\end{minipage}\par\vfill
\noindent\begin{minipage}{\textwidth}\ProblemHeading{习题}{3.182}
设 $f$ 在区间 $I$ 上连续，且任意 $x_1,x_2\in I$ 都满足\[f\left(\frac{x_1+x_2}{2}\right)\le\frac{f(x_1)+f(x_2)}2.\]证明 $f$ 是区间 $I$ 上的凸函数。
\end{minipage}\par\vfill
\noindent\begin{minipage}{\textwidth}\ProblemHeading{习题}{3.183}
设 $f$ 是 $[0,1]$ 上的二阶可微函数，且 $f(0)=f(1)=0$。证明存在 $\xi\in(0,1)$，使\[f^{\prime\prime}(\xi)=\frac{2f^{\prime}(\xi)}{1-\xi}.\]
\end{minipage}\par\vfill
\noindent\begin{minipage}{\textwidth}\ProblemHeading{习题}{3.184}
设 $f$ 在 $[a,b]$ 上连续，且 $f(a)=f(b)=0$，端点的单侧导数满足 $f^{\prime}(a)f^{\prime}(b)>0$。证明存在 $\xi\in(a,b)$，使 $f(\xi)=0$。
\end{minipage}\par\vfill
\newpage
\noindent\begin{minipage}{\textwidth}\ProblemHeading{习题}{3.185}
设 $f$ 在 $[0,1]$ 上可导，$f(0)=1$、$f(1)=1/2$。证明存在 $\xi\in(0,1)$，使 $f^2(\xi)+f^{\prime}(\xi)=0$。
\end{minipage}\par\vfill
\noindent\begin{minipage}{\textwidth}\ProblemHeading{习题}{3.186}
设 $f$ 在 $[a,+\infty)$ 上二阶可微，$f(a)>0$、$f^{\prime}(a)<0$，并且对 $x>a$ 有 $f^{\prime\prime}(x)\le0$。证明 $f$ 在 $(a,+\infty)$ 内恰有一个零点。
\end{minipage}\par\vfill
\noindent\begin{minipage}{\textwidth}\ProblemHeading{习题}{3.187}
设 $f$ 在 $[a,b]$ 上可微，$f^{\prime}$ 严格递增，且 $f(a)=f(b)=\lambda$。证明对任意 $x\in(a,b)$ 都有 $f(x)<\lambda$。
\end{minipage}\par\vfill
\noindent\begin{minipage}{\textwidth}\ProblemHeading{习题}{3.188}
设 $f$ 在 $(a,+\infty)$ 上可导，且 $\lim_{x\to+\infty}f(x)$ 为有限数。
\SubQ{1}{用 $f(x)=\sin(x^2)/x\ (x>0)$ 说明，不能保证 $\lim_{x\to+\infty}f^{\prime}(x)$ 存在。}
\SubQ{2}{证明：若 $\lim_{x\to+\infty}f^{\prime}(x)$ 存在，则其值必为零。}
\end{minipage}\par\vfill
\newpage
\noindent\begin{minipage}{\textwidth}\ProblemHeading{习题}{3.189（校订）}
设 $f$ 在 $[0,+\infty)$ 上连续，在 $(0,+\infty)$ 内二阶可微，$f^{\prime\prime}(x)<0$ 且 $f(0)=0$。证明对任意正数 $x_1,x_2$ 都有 $f(x_1+x_2)<f(x_1)+f(x_2)$。
\end{minipage}\par\vfill
\noindent\begin{minipage}{\textwidth}\ProblemHeading{习题}{3.190}
设 $f$ 在 $x_0$ 处存在二阶导数。求\[\lim_{h\to0}\frac{f(x_0+h)+f(x_0-h)-2f(x_0)}{h^2}.\]
\end{minipage}\par\vfill
\noindent\begin{minipage}{\textwidth}\ProblemHeading{习题}{3.191}
证明下列不等式。
\SubQ{1}{对任意实数 $x$，$e^x\ge1+x+\dfrac{x^2}{2}+\dfrac{x^3}{6}$。}
\SubQ{2}{对 $x>0$，$x-\dfrac{x^2}{2}\le\ln(1+x)\le x-\dfrac{x^2}{2}+\dfrac{x^3}{3}$。}
\SubQ{3}{对 $0<x<\pi/2$，$x-\dfrac{x^3}{6}<\sin x<x-\dfrac{x^3}{6}+\dfrac{x^5}{120}$。}
\end{minipage}\par\vfill
\noindent\begin{minipage}{\textwidth}\ProblemHeading{习题}{3.192}
证明下列不等式。
\SubQ{1}{对任意实数 $x,y$，$2e^{(x+y)/2}\le e^x+e^y$。}
\end{minipage}\par\vfill
\newpage
\noindent\begin{minipage}{\textwidth}\ProblemHeading{习题}{3.193}
求 $\displaystyle\lim_{n\to\infty}\left(1+\frac1{n^2}\right)\left(1+\frac2{n^2}\right)\cdots\left(1+\frac n{n^2}\right)$。
\end{minipage}\par\vfill
\noindent\begin{minipage}{\textwidth}\ProblemHeading{习题}{3.194}
求数列 $\sqrt[n]{n}\ (n=1,2,\ldots)$ 的最大值。
\end{minipage}\par\vfill
\noindent\begin{minipage}{\textwidth}\ProblemHeading{习题}{3.195}
给出函数 $x\cos x$ 在 $[0,\pi/2]$ 上的一个尽可能小的上界。
\end{minipage}\par\vfill
\noindent\begin{minipage}{\textwidth}\ProblemHeading{习题}{3.196}
设 $f$ 在 $[-1,1]$ 上有三阶连续导数，且 $f(-1)=0$、$f(1)=1$、$f^{\prime}(0)=0$。证明存在 $\xi\in(-1,1)$，使 $f^{(3)}(\xi)=3$。
\end{minipage}\par\vfill
\newpage
\noindent\begin{minipage}{\textwidth}\ProblemHeading{习题}{3.197}
设 $a>1$，函数 $f:(0,+\infty)\to(0,+\infty)$ 可微。证明存在趋于正无穷的正数列 $\{x_n\}$，使 $f^{\prime}(x_n)<f(ax_n)$ 对所有 $n=1,2,\ldots$ 成立。
\end{minipage}\par\vfill
\noindent\begin{minipage}{\textwidth}\ProblemHeading{习题}{3.198}
利用凸函数的性质证明 Hölder 不等式：设 $a_i,b_i>0\ (i=1,\ldots,n)$，$p,q>1$ 且 $1/p+1/q=1$，则\[\sum_{i=1}^na_ib_i\le\left(\sum_{i=1}^na_i^p\right)^{1/p}\left(\sum_{i=1}^nb_i^q\right)^{1/q}.\]
\end{minipage}\par\vfill
